% Requisitos funcionales del sistema (IT-59).
%
% La ultima columna no es decoracion: separa los requisitos que salen de una
% decision de diseno (remiten a una seccion) de los que salen de un experimento
% (remiten a un ADR). Esa distincion es la que pide la tarjeta.
%
% Las columnas se reparten el ancho con \hsize porque el enunciado del
% requisito necesita mas sitio que la procedencia; los dos factores tienen que
% sumar 2 (el numero de columnas X).
\begin{table}[htbp]
  \centering
  \small
  \caption[Requisitos funcionales del sistema]{Requisitos funcionales del
    sistema. Fuente: Elaboración propia.}
  \label{tab:requisitos_funcionales}
  \begin{tabularx}{\textwidth}{@{}l >{\hsize=1.42\hsize}Y c >{\hsize=0.58\hsize}Y@{}}
    \toprule
    \textbf{ID} & \textbf{Requisito} & \textbf{Traza} & \textbf{Se fija en} \\
    \midrule
    RF-01 & Construir la colección documental a partir de la web oficial de la EPSJ, conservando la procedencia y fecha de recolección de cada documento                & OE1 · RU-02 & ADR-0002 \\
    RF-02 & Fragmentar la colección sin que ningún fragmento mezcle dos asignaturas                                                          & OE2 · RU-02 & ADR-0001 \\
    RF-03 & Indexar los fragmentos en una base de datos vectorial que se reconstruya por completo en cada ejecución                          & OE3         & ADR-0004 \\
    RF-04 & Recuperar de esa base los fragmentos pertinentes a cada consulta, en número variable según su pertinencia                        & OE4 · RU-01 & Sección~\ref{secc:fase2_pipeline} \\
    RF-05 & Redactar la respuesta empleando únicamente los fragmentos recuperados y no información inventada                                 & OE4 · RU-02 & ADR-0005 \\
    RF-06 & Retirar la respuesta entera, y no solo el nombre, cuando nombre una titulación que no está en el catálogo                        & OE4 · RU-03 & Sección~\ref{secc:fase2_dominio} \\
    RF-07 & Avisar de que la pregunta queda fuera de su ámbito en lugar de improvisar una respuesta                                          & OE4 · RU-06 & Sección~\ref{secc:fase2_dominio} \\
    RF-08 & Resolver las preguntas de seguimiento manteniendo el hilo de los últimos turnos de la conversación                              & OE6 · RU-04 & Sección~\ref{secc:fase2_conversacion} \\
    RF-09 & Mostrar junto a la respuesta las fuentes de la colección de las que sale                                                         & OE5 · RU-02 & Sección~\ref{secc:fase2_pipeline} \\
    RF-10 & Ofrecer la interfaz de chat sin registro previo ni instalación por parte del usuario                                             & OE6 · RU-05 & Sección~\ref{secc:fase3_interfaz} \\
    \bottomrule
  \end{tabularx}
  \begin{minipage}{0.95\textwidth}
    \vspace{4pt}
    \tiny
    \textbf{Notas:} \textbf{RF} identifica un requisito funcional del sistema.
    La columna \textbf{Traza} enlaza cada uno con el objetivo específico
    (\textbf{OE}) que lo respalda y con el requisito de usuario (\textbf{RU})
    del que nace, recogido en la Tabla~\ref{tab:requisitos_usuario}. La columna
    \textbf{Se fija en} distingue el origen del requisito: los que remiten a un
    \textbf{ADR} (\textit{Architecture Decision Record}) dependen de una
    decisión tomada por experimentación, los que remiten a una sección salen de una decisión
    de diseño argumentada en el texto.
  \end{minipage}
\end{table}
